\documentclass{beamer}
\usepackage{beamerthemeshadow}
\usepackage{amsmath,amsfonts,amsthm}

\setbeamertemplate{navigation symbols}{}
\usetheme{Warsaw}
\beamersetuncovermixins{\opaqueness<1>{25}}{\opaqueness<2->{15}}

\include{Qcircuit}
\newcommand{\E}{\mathbf{E}}
\newcommand{\cA}{\mathcal A}
\newcommand{\cB}{\mathcal B}
\newcommand{\cH}{\mathcal H}
\newcommand{\C}{\mathbb C}
\newcommand{\TVD}{\mathrm{TVD}}
\renewcommand{\>}{\rangle}
\newcommand{\<}{\langle}


\begin{document}
\title{Algorithm for Preparing Thermal Gibbs States}  
\author{Pawel Wocjan\\ University of Central Florida}
\date{\today} 
\frame{\titlepage}

%%%%%%%%%%%%%%%%%%%%%%%%%%%%%%%%%%%%%%%%%%%%%%%%%%%%%%%%%%%%%%%%%%%%%%%%%%%%
\frame{\frametitle{Related work}

This work is based on

\medskip

\begin{itemize}
\item {\em Sampling from the Thermal Quantum Gibbs State and Evaluating Partition Functions with a Quantum Computer} 

\medskip
with David Poulin, Phys. Rev. Lett. 103, 220502 (2009)

\bigskip	
\item {\em Algorithm for Preparing Thermal States -- Detailed Analysis} 

\medskip
with C.-F. Chiang, Proc. NATO ARW ``Quantum Cryptography and Computing: Theory and Implementation, Gdansk, Poland, 9-12 September 2009  
\end{itemize}

\bigskip
This work has been supported by the National Science Foundation grants CCF-0726771 and CCF-0746600
} 

%%%%%%%%%%%%%%%%%%%%%%%%%%%%%%%%%%%%%%%%%%%%%%%%%%%%%%%%%%%%%%%%%%%%%%%%%%%%%%%%%%%%%%%%%%%

\frame{\frametitle{Definitions}

Let $H$ be a Hamiltonian with spectral decomposition
\[
H = \sum_{a=1}^D E_a |\psi_a\>\<\psi_a|\ %,, \quad E_a \in [0,1]
\]

The partition function $Z_\beta$ at inverse temperature $\beta$ is
\[
Z_\beta := \sum_{a=1}^D e^{-\beta E_a}
\]

The thermal state $\rho_\beta$ at inverse temperature $\beta$ is
\[
\rho_\beta := \sum_{a=1}^D \frac{e^{-\beta E_a}}{Z_\beta} |\psi_a\>\<\psi_a|
\]
}

%%%%%%%%%%%%%%%%%%%%%%%%%%%%%%%%%%%%%%%%%%%%%%%%%%%%%%%%%%%%%%%%%%%%%%%%%%%%%%%

\frame{\frametitle{Thermalizing quantum systems}

Given $H$, $\beta$, and $\epsilon$, prepare a quantum state $\tilde{\rho}_\beta$ with

\medskip
\[
\| \rho_\beta - \tilde{\rho}_\beta \|_{\mathrm{tr}} \le \epsilon 
\]

\bigskip

}

%%%%%%%%%%%%%%%%%%%%%%%%%%%%%%%%%%%%%%%%%%%%%%%%%%%%%%%%%%%%%%%%%%%%%%%%%%%%%%%%%%%%%%%%%%%%%%%%%

\frame{\frametitle{Methods}
Metropolis Sampling
  
\bigskip

Quantized Metropolis Sampling  

\begin{itemize}
\item Somma, Boixo, Barnum, and Knill, PRL 101,130504 
\item W. and Abeyesinghe, PRA 78, 032311
\end{itemize}
  
\bigskip

Method based on Amplitude Amplification (this presentation)
  
\bigskip

Quantum Metropolis Sampling (David Poulin's presentation)

}

%%%%%%%%%%%%%%%%%%%%%%%%%%%%%%%%%%%%%%%%%%%%%%%%%%%%%%%%%%%%%%%%%%%%%%%%%%%%%%%%%%%%%%%%%%%%%%%%%%

\frame{\frametitle{Classical Hamiltonians}

Let $\Omega$ be a set whose elements $x$ correspond to the states of some classical system

\bigskip
Let $E : \Omega\rightarrow [0,1]$ denote the energy function, assigning to each state $x$ its energy $E(x)$

\bigskip
Let 
\[
H = \sum_{x\in\Omega} E(x) |x\>\<x|
\]

\bigskip 
Thermalizing such classical systems is easier because $H$ is diagonal in the computational basis $\{ |x\> \}_{x\in\Omega}$

}

%%%%%%%%%%%%%%%%%%%%%%%%%%%%%%%%%%%%%%%%%%%%%%%%%%%%%%%%%%%%%%%%%%%%%%%%%%%%%%%%%%%%%%%%%%%%%%%%%%%%%%

\frame{\frametitle{Metropolis sampling}

Given such classical Hamiltonian $H$, we can always construct a Markov chain $P$ whose limiting distribution is equal to $\rho_\beta$

\bigskip
The spectral gap $\delta$ determines how fast we can approach the Boltzmann distribution

\bigskip
The run time scales like $1/\delta$

\bigskip
In general, it is not possible to determine (or bound) the spectral gap $\Rightarrow$ heuristics

\bigskip
There are important cases where it can be shown that the gap is large

}

%%%%%%%%%%%%%%%%%%%%%%%%%%%%%%%%%%%%%%%%%%%%%%%%%%%%%%%%%%%%%%%%%%%%%%%%%%%%%%%%%%%%%%%%%%%%%%%%%%%%%%%%%%%%%%%

\frame{\frametitle{Quantized Metropolis algorithm}

Given such classical Hamiltonian $H$, there is a ``quantized'' Metropolis
algorithm preparing a state $|\tilde{\varphi}_\beta\>$ with

\medskip
\[
\big\| |\varphi_\beta\> - |\tilde{\varphi}_\beta\> \big\| \le \epsilon
\]

\bigskip
where $|\varphi_\beta\>$ is a ``quantized'' version of the Boltzmann distribution

\medskip
\[
|\varphi_\beta\> := \sum_{x\in \Omega} \sqrt{\frac{e^{-\beta E_x}}{Z_\beta}} |x\>
\]

\bigskip
The run times scales like $\sqrt{1/\delta}$ 

}

%%%%%%%%%%%%%%%%%%%%%%%%%%%%%%%%%%%%%%%%%%%%%%%%%%%%%%%%%%%%%%%%%%%%%%%%%%%%%%

\frame{\frametitle{Structure of our method}

Prepare the state 
\[
\frac{1}{\sqrt{D}}\sum_{x\in\Omega} |x\> \otimes |E_x\> \otimes |0\>
\]

\bigskip
Apply a controlled rotation to obtain the state
\[
|\Psi\> = \frac{1}{\sqrt{D}}\sum_{x\in\Omega} |x\> \otimes |E_x\> \otimes \Big(
\sqrt{e^{-\beta E_x}} |0\> + \sqrt{1-e^{-\beta E_x}} |1\> \Big)
\]

}

%%%%%%%%%%%%%%%%%%%%%%%%%%%%%%%%%%%%%%%%%%%%%%%%%%%%%%%%%%%%%%%%%%%%%%%%%%%%%%

\frame{\frametitle{Structure of our method}

Measure in the basis $|0\>$, $|1\>$

\bigskip
If we obtain $|0\>$ $\Rightarrow$ the post-measurement state is

\medskip
\[
\frac{|\Pi_0 |\Psi\>}{\|\Pi_0 |\Psi\>\|} = 
\sum_{x\in\Omega} \sqrt{\frac{e^{-\beta E_x}}{Z_\beta}} |x\> \otimes |E_x\> \otimes  |0\> 
\]

\medskip
where $\Pi_0 = I \otimes |0\>\<0|$

\bigskip
Using amplitude amplification, we can achieve an expected run time
\[
\frac{1}{\|\Pi_0 |\Psi\>\|} = \sqrt{\frac{D}{Z_\beta}}
\]

}

%%%%%%%%%%%%%%%%%%%%%%%%%%%%%%%%%%%%%%%%%%%%%%%%%%%%%%%%%%%%%%%%%%%%%%%%%%%%%%%%%%

\frame{\frametitle{Generalization to quantum Hamiltonians}

Let $H$ be an arbitrary quantum Hamiltonian with spectral decomposition

\medskip
\[
H = \sum_{a=1}^D E_a |\psi_a\>\<\psi_a| \quad\quad E_a \in [0,1/2]
\]

\bigskip
There are two difficulties:

\medskip
\begin{itemize}
\item we cannot change between the computational basis $|a\>$ and the eigenvector basis $|\psi_a\>$ of $H$

\medskip
\item we do not know the energies $E_a$
\end{itemize}

\bigskip
Both difficulties can be overcome with quantum phase estimation $\Rightarrow$ we can extend the previous method to quantum Hamiltonians
}

%%%%%%%%%%%%%%%%%%%%%%%%%%%%%%%%%%%%%%%%%%%%%%%%%%%%%%%%%%%%%%%%%%%%%%%%%%%%%%%%%%

\frame{\frametitle{Error sources}

Given $H$, $\beta$, and $\epsilon$, our quantum algorithm prepares $\tilde{\rho}_\beta$
\[
 \| \rho_\beta - \tilde{\rho}_\beta \|_{\mathrm{tr}} \le \epsilon 
\]

\bigskip
The errors come from two sources:
\[
\| \rho_\beta           - \tilde{\rho}_\beta   \|_{\mathrm{tr}} \le 
\| \rho                 - \rho_{\textrm{sim}}  \|_{\mathrm{tr}} +
\| \rho_{\textrm{sim}}  - \rho_{\textrm{fail}} \|_{\mathrm{tr}} 
\]

\bigskip
	\begin{itemize}
	\item $\rho_{\textrm{sim}}$ due to imperfect simulation of $\exp(2 \pi i H)$
	
	\bigskip
	\item $\tilde{\rho}_\beta=\rho_{\textrm{fail}}$ due to nonzero failure
	probability and limited precision of quantum phase estimation
	\end{itemize}

}

%%%%%%%%%%%%%%%%%%%%%%%%%%%%%%%%%%%%%%%%%%%%%%%%%%%%%%%%%%%%%%%%%%%%%%%%%%%%%%%%

\frame{\frametitle{Imperfect simulation of Hamiltonians}

Assume that $H$ is sparse and let $U = \exp(2\pi i H)$

\bigskip
Using techniques for simulating Hamiltonian time evolutions, we can implement $U_{\textrm{sim}}$ with

\medskip
\[
\| U - U_{\textrm{sim}} \| \le \epsilon_{\textrm{sim}}
\]

\bigskip
The cost scales like 
\[
\mathrm{poly}(\log D, 1/\epsilon_{\textrm{sim}}, s)
\]

\medskip
where $s$ indicates the sparseness of $H$

}

%%%%%%%%%%%%%%%%%%%%%%%%%%%%%%%%%%%%%%%%%%%%%%%%%%%%%%%%%%%%%%%%%%%%%%%%%%%%%%%%

\frame{\frametitle{Technical results}

Matrix logarithm is Lipschitz continuous:

\[
\| U - U_{\textrm{sim}} \| \le \epsilon_{\textrm{sim}}
\]

$\Rightarrow$ there is a Hamiltonian $H_{\textrm{sim}}$ with
\[
\| H - H_{\textrm{sim}} \| \le O(\epsilon_{\textrm{sim}})
\]

\bigskip
Thermal states are H\"older continuous:

\[
\| H - H_{\textrm{sim}} \| \le O(\epsilon_{\textrm{sim}})
\]

$\Rightarrow$ the corresponding thermal states satisfy
\[
\| \rho - \rho_{\textrm{sim}} \|_{\textrm{tr}} \le O(\sqrt{\beta \epsilon_{\textrm{sim}}})
\]    
 
}

%%%%%%%%%%%%%%%%%%%%%%%%%%%%%%%%%%%%%%%%%%%%%%%%%%%%%%%%%%%%%%%%%%%%%%%%%%%%%%%%

\frame{\frametitle{Structure of our method}

Start with the maximally entangled state 

\medskip
\[
|\Phi\>:=\frac{1}{\sqrt{D}} \sum_{a=1}^D |a\> \otimes |a\>
\]

\bigskip
Run phase estimation of $U_\mathrm{sim}$ on first part of $|\Phi\>$ and record eigenphase
in the energy register

\bigskip
Apply the controlled rotation $R = \sum_E |E\>\<E| \otimes R_E$	 

\medskip
\[
R_E =  
\left(
\begin{array}{cc}
\sqrt{e^{-\beta E}}   & - \sqrt{1-e^{-\beta E}} \\
\sqrt{1-e^{-\beta E}} &   \sqrt{e^{-\beta E}}
\end{array}
\right)
\] 

\bigskip
Use amplitude amplification to increase overlap with $|0\>$ component
}

%%%%%%%%%%%%%%%%%%%%%%%%%%%%%%%%%%%%%%%%%%%%%%%%%%%%%%%%%%%%%%%%%%%%%%%%%%%%%%%%%

\frame{\frametitle{Invariance of maximally entangled state}

Let
\begin{eqnarray*}
|\Phi\>  & := & \frac{1}{\sqrt{D}} \sum_{a=1}^D |a\> \otimes |a\> \\ & & \\
V        & := & \sum_{a=1}^D |\psi_a\>\<a|
\end{eqnarray*}

\bigskip
Due to invariance of $|\Phi\>$ under $V\otimes \bar{V}$, we have

\medskip
\[
|\Phi\> = \frac{1}{\sqrt{D}} \sum_{a=1}^D |\psi_a\> \otimes |\bar{\psi}_a\>
\]

}

%%%%%%%%%%%%%%%%%%%%%%%%%%%%%%%%%%%%%%%%%%%%%%%%%%%%%%%%%%%%%%%%%%%%%%%%%%%%%%%%%

\frame{\frametitle{Correspondance between energy levels and eigenphases}

Recall that $U=\exp(2\pi iH)$

\bigskip
We have
\begin{eqnarray*}
(U \otimes I_D) |\psi_a\> \otimes |\bar{\psi}_a\> & = & e^{2 \pi i E_a}
|\psi_a\> \otimes |\bar{\psi}_a\> \\ && \\
(U \otimes I_D) |\Phi\>                       & = & \frac{1}{\sqrt{D}} \sum_{a=1}^D e^{2 \pi i E_a}
|\psi_a\> \otimes |\bar{\psi}_a\>
\end{eqnarray*}

The energy levels $E_a$ of $H$ correspond bijectively to the eigenphases
$E_a$ of $U$

\bigskip
$\Rightarrow$ Apply quantum phase estimation to approximate the energy levels

}

%%%%%%%%%%%%%%%%%%%%%%%%%%%%%%%%%%%%%%%%%%%%%%%%%%%%%%%%%%%%%%%%%%%%%%%%%%%%%%%%%%

\frame{\frametitle{Intuition -- perfect quantum phase estimation}

If perfect quantum phase estimation were possible (infinite precion and zero probability of failure), we could prepare the state 
\[
|\Psi\> = \frac{1}{\sqrt{D}}
\sum_{a=1}^D |\psi_a\> \otimes |\bar{\psi}_a\> \otimes |E_a\> 
\otimes 
\Big(
\sqrt{e^{-\beta E_a}}  |0\> + \sqrt{1-e^{-\beta E_a}} |1\> 
\Big)
\]

\bigskip
$\Rightarrow$ Using amplitude amplification just as in the case of classical Hamiltonians,
we could obtain a purification of the thermal state

}

%%%%%%%%%%%%%%%%%%%%%%%%%%%%%%%%%%%%%%%%%%%%%%%%%%%%%%%%%%%%%%%%%%%%%%%%%

\frame{\frametitle{Imperfect quantum phase estimation}

We use quantum phase estimation that invokes the controlled $U_\mathrm{sim}$ operation

\[
O((1/\epsilon_{\textrm{prec}})\log(1/\epsilon_{\textrm{fail}})) \mbox{ times} 
\]

\bigskip
$\Rightarrow$ We transform $|\Phi\>$ to the state 
\[
\frac{1}{\sqrt{D}} 
\sum_{a=1}^D |\psi_a\> \otimes |\bar{\psi}_a\> \otimes \Big( c_a^- |E_a^-\> +
c_a^+ |E_a^+\> \Big) + |\xi\>
\]

\medskip
\begin{itemize}
\item $E_a^+ - E_a  \le \epsilon_{\textrm{prec}}$, \quad $E_a - E_a^- \le \epsilon_{\textrm{prec}}$

\medskip
\item $|c_a^+|^2 + |c_a^-|^2 \geq 1 - \epsilon_{\textrm{fail}}$

\medskip
\item $\| |\xi\> \| ^2 \le \epsilon_{\textrm{fail}}$
\end{itemize}

}

%%%%%%%%%%%%%%%%%%%%%%%%%%%%%%%%%%%%%%%%%%%%%%%%%%%%%%%%%%%%%%%%%%%%%%%%%

\frame{\frametitle{Controlled rotation}

After quantum phase estimation, we apply the rotation controlled by the value in the energy register to prepare the state $|\tilde{\Psi}\>$ 

\medskip
\[ 
\frac{1}{\sqrt{D}} \sum_{a=1}^D |\psi_a\> \otimes |\bar{\psi}_a\> \otimes c_a^\pm |E_a^\pm\> \otimes 
\Big(
\sqrt{e^{-\beta E_a^\pm}} |0\> + \sqrt{1-e^{-\beta E_a^\pm}} |1\>
\Big) 
 + |\tilde{\xi}\>
\]

}

%%%%%%%%%%%%%%%%%%%%%%%%%%%%%%%%%%%%%%%%%%%%%%%%%%%%%%%%%%%%%%%%%%%%%%%%%%%%%%%%%%

\frame{\frametitle{Amplitude amplification}

We apply amplitude amplification to obtain

\medskip
\[
|\tilde{\beta}\> := \frac{\Pi_0 |\tilde{\Psi}\>}{\|\Pi_0 |\tilde{\Psi}\>\|} 
\] 

\bigskip
The required number of iterations is 

\medskip
\[
1/\|\Pi_0 |\tilde{\Psi}\>\|
\]

\bigskip
The desired approximate thermal state is 

\medskip
\[\tilde{\rho}_\beta =
\mathrm{Tr}_{\bar{A}}(|\tilde{\beta}\>\<\tilde{\beta}|)
\]
where we trace out over the ancilla, the energy register and the register holding ${|\bar{\psi}_a}\>$ 

}

%%%%%%%%%%%%%%%%%%%%%%%%%%%%%%%%%%%%%%%%%%%%%%%%%%%%%%%%%%%%%%%%%%%%%%%%%

\frame{\frametitle{Error analysis}

The approximate thermal state $\tilde{\rho}_\beta$ is close to the perfect $\rho_\beta$ provided that simulation and phase estimation are accurate enough

\bigskip
Given $H$, $\beta$, and $\epsilon$, our quantum algorithm makes it possible to prepare $\tilde{\rho}_\beta$

\medskip
\[
\| \rho_\beta - \tilde{\rho}_\beta \|_{\mathrm{tr}} \le \epsilon 
\]

\bigskip
The expected run time is

\medskip
\[
\sqrt{\frac{D}{Z_\beta}} \cdot \frac{\beta}{\epsilon} \cdot \left(
\log \frac{1}{\epsilon} +
\beta
\right)
\]

}
%%%%%%%%%%%%%%%%%%%%%%%%%%%%%%%%%%%%%%%%%%%%%%%%%%%%%%%%%%%%%%%%%%%%%%%%%%
\end{document}

